% Part: normal-modal-logic
% Chapter: syntax-and-semantics
% Section: normal-modal-logics

\documentclass[../../../include/open-logic-section]{subfiles}

\begin{document}

\olfileid[fr]{nml}{syn}{nor}

\olsection{Logiques modales normales}

On peut concevoir une logique modale comme un ensemble de !!{formula}s
qui satisfait certaines conditions de fermeture. Quelques logiques
modales présentent un intérêt particulier par leurs interprétations
et leurs applications, mais surtout parce que leurs propriétés
théoriques sont bien connues. Ce sont les logiques modales dites normales.

\begin{defn}\ollabel{defn:modal-logic}
Une \emph{logique modale}~$\Log L$ est un ensemble de !!{formula}s
du langage modal de base qui
\begin{enumerate}
\item contient toutes les tautologies propositionnelles;
\item est clos par substitution uniforme;
\item est clos par modus ponens : si $!A \in \Log L$ et
  $(!A \lif !B) \in \Log L$, alors $!B \in \Log L$.
\end{enumerate}
\end{defn}

\begin{ex}
Les ensembles suivants satisfont cette définition et constituent
donc des logiques modales, bien qu'ils présentent peu d'intérêt :
\begin{enumerate}
\item l'ensemble de toutes les tautologies;
\item l'ensemble de toutes les !!{formula}s, c'est-à-dire la logique
  modale incohérente.
\end{enumerate}
\end{ex}

\begin{defn}\label{defn:normal-modal-logic}
Une logique modale~$\Log L$ est \emph{normale} si et seulement si
\begin{enumerate}
\item elle contient les deux !!{formula}s
\begin{align}
\tag{K} \Box(p \lif q) & \lif (\Box p \lif \Box q) \\
\tag{Dual} \Diamond p & \liff \lnot \Box \lnot p
\end{align}
\item elle est close par nécessitation : si $!A \in
  \Log L$, alors $\Box !A \in \Log L$.
\end{enumerate}
\end{defn}

Nous rencontrerons deux façons principales de définir une logique
modale. La théorie de la démonstration la décrit comme l'ensemble
des !!{formula}s !!{derivable}s dans un certain système de
!!{derivation}. La théorie des modèles la décrit comme l'ensemble
des !!{formula}s valides dans une certaine classe de modèles.
On retrouve ainsi la situation de la logique propositionnelle
ordinaire, et aussi de la logique du premier ordre : on peut
caractériser le même ensemble de !!{formula}s comme l'ensemble des
tautologies, vraies sous toute assignation de valeurs de vérité,
ou comme l'ensemble des !!{formula}s !!{derivable}s dans un système
de !!{derivation}. Pour les logiques modales normales, ces deux
approches font appel aux modèles relationnels et aux systèmes
axiomatiques de démonstration de type Hilbert.

\end{document}
